\documentclass[11pt]{article}

% --- packages ---
\usepackage[margin=1in]{geometry}
\usepackage{amsmath,amssymb,amsfonts}
\usepackage{booktabs}
\usepackage{graphicx}
\usepackage{hyperref}
\usepackage{xcolor}
\usepackage{multirow}
\usepackage{float}
\usepackage[numbers]{natbib}
\usepackage{algorithm}
\usepackage{algorithmic}
% \usepackage{microtype}  % disabled: MiKTeX font expansion issue
\usepackage{bookmark}

\hypersetup{
  colorlinks=true,
  linkcolor=blue!60!black,
  citecolor=blue!60!black,
  urlcolor=blue!60!black
}

% --- metadata ---
\title{%
  The Topology Programmer:\\
  Cybernetic Feedback as a General-Purpose Geometric\\
  Prior for Neural Network Adapters%
}

\author{%
  Timothy Paul Bielec%
  \thanks{Corresponding author. \texttt{timothy@promptcrafted.com}}%
  \\
  Promptcrafted LLC\\
  Los Angeles, CA
}

\date{March 2026}

\begin{document}
\maketitle

% ===================================================================
\begin{abstract}
A companion paper established that cybernetic feedback drives
block-diagonal emergence in 6-channel LoRA bridges aligned with
rhombic dodecahedral face-pair geometry (co/cross 70{,}404:1). We show
the mechanism is general across four polytope geometries: octahedron
($2{+}2$ blocks, $n{=}4$, co/cross 473{,}622:1), rhombic dodecahedron
($3{+}3$ blocks, $n{=}6$), tesseract ($4{+}4$ blocks, $n{=}8$, co/cross
41{,}564:1), and 24-cell ($12{+}12$ blocks, $n{=}24$, co/cross
35{,}808:1)---with identical Steersman, identical hyperparameters, and
only the pair specification changed. Channel ablation across
$n \in \{3, 4, 6, 8, 12\}$ reveals a universal spectral attractor at
Fiedler $0.084$--$0.102$ for spectral-only training: the bridge's
natural connectivity equilibrium. Contrastive loss breaks this
equilibrium by $1{,}130\times$ at $n{=}6$. Two negative controls
establish the geometric coherence requirement: wrong-labels (random
3-pair partition) produces clean $3{+}3$ eigenvalue splits but co/cross
${\sim}\,10^{-5}$:1---eigenvalue structure without directional content.
Resonance (number-theoretic pairs) produces chain-like eigenvalues with
no block structure. An emanation architecture (master bridge with
per-layer offsets) converges to the spectral attractor, confirming
hierarchy provides connectivity without directionality. These outcomes
define four cleanly separated regimes: block-diagonal, spectral
attractor, wrong-partition, and collapse. The bridge is not a
geometry-discovering mechanism---it is a programmable topology substrate
that accepts only geometrically coherent programs. A Steersman annealing
experiment reveals five control regimes: the optimal operating point is
at ${\sim}2\%$ of nominal weight (co/cross 18{,}671:1), not full
strength, which operates in a destructive interference regime. A
whisper-strength experiment ($c_w = 0.02$, adaptive decay to floor
$0.005$) reveals a three-phase trajectory: rapid growth, incubation
plateau at ${\sim}1{,}700$:1, then exponential breakout to peak
15{,}183:1, stabilizing at 13{,}456:1---$5.2\times$ below the
aggressive regime, suggesting $c_w$ controls at minimum the rate and
possibly the ceiling of structural formation. Training is reproducible:
tesseract contrastive at $n{=}8$ replicates with Pearson $r = 1.0000$
across independent runs.
\end{abstract}

\noindent\textbf{Keywords:} topology programming, cybernetic feedback,
bridge matrix, polytope geometry, spectral attractor, multi-channel
LoRA, parameter-efficient fine-tuning, block-diagonal structure


% ===================================================================
\section{Introduction}
\label{sec:intro}

Low-rank adaptation~\cite{hu2022lora} modifies pretrained language models
by injecting rank-$r$ update matrices $\Delta W = BA$ alongside frozen
weights. Multi-channel extensions~\cite{bielec2026weights,bielec2026bridge}
partition the rank dimension into $n$ channels governed by an
$n \times n$ learnable bridge matrix that controls inter-channel coupling.
The bridge adds $n^2$ parameters (36 at $n{=}6$) and introduces a degree
of freedom absent from standard LoRA: the \emph{topology} of the adapter.

Paper~3~\cite{bielec2026bridge} demonstrated that a closed-loop cybernetic
controller---the Steersman---drives 6-channel bridges toward 3-axis
block-diagonal structure aligned with the rhombic dodecahedron's (RD)
face-pair geometry. The result was consistent across six experiments,
42{,}500+ bridges, two model scales, and three initialization strategies,
with coupling ratios exceeding 22{,}000:1. Paper~3 framed this as
geometry \emph{discovery}: the Steersman finds structure latent in the
loss landscape.

This paper reframes the result. The bridge is not a discovery mechanism
but a \emph{programmable substrate}. The Steersman is not a geometer but
a \emph{topology compiler}. The pair specification---the declaration of
which channels are co-axial---is a program. Paper~3 ran one program (RD
face-pairs at $n{=}6$) and observed one output (three $2 \times 2$
blocks). The question this paper answers: does the compiler accept other
programs, and what happens when it receives an invalid one?

The answer is that the Steersman is general-purpose but selective. It
compiles any pair specification derived from polytope co-axial structure
and rejects specifications that lack genuine geometric coherence. With no
modification to the feedback mechanism, control laws, or training
hyperparameters---only the pair specification changes---the Steersman
programs three polytope geometries across two spatial dimensionalities:

\begin{itemize}
  \item \textbf{Octahedron} ($n{=}4$, 3D, 2 co-axial pairs):
    two $2 \times 2$ blocks, per-bridge mean coupling ratio
    \textbf{473{,}622:1}---the programme's strongest signal.
  \item \textbf{Rhombic dodecahedron} ($n{=}6$, 3D, 3 co-planar pairs):
    three $2 \times 2$ blocks, coupling ratio \textbf{70{,}404:1}.
    Paper~3's result, reproduced across seeds.
  \item \textbf{Tesseract} ($n{=}8$, 4D, 4 co-axial pairs):
    four $2 \times 2$ blocks, coupling ratio \textbf{41{,}564:1},
    Fiedler eigenvalue 0.000191, clean $4{+}4$ eigenvalue split.
\end{itemize}

\noindent Four negative controls define the compiler's rejection
boundary. Spectral-only training across $n \in \{3, 4, 8, 12\}$ reveals
a universal spectral attractor at Fiedler ${\approx}\,0.09$---the
bridge's natural connectivity equilibrium when feedback provides no
topological direction. Wrong-labels (WL-001: random 3-pair partition on
$n{=}6$) produces the predicted $3{+}3$ eigenvalue split but co/cross
ratio ${\sim}\,10^{-5}$:1---eigenvalue structure without directional
content. Resonance (R-001: number-theoretic pairs) produces chain-like
eigenvalues rather than block structure. An emanation architecture
(E-001: hierarchical master bridge with per-layer offsets) converges to
the spectral attractor, not to block-diagonal structure---hierarchy
provides connectivity without directionality.

These outcomes define four cleanly separated dynamical regimes with no
intermediate cases: block-diagonal ($\rho > 40{,}000$:1), spectral
attractor (Fiedler~${\approx}\,0.09$, $\rho \approx 1$:1),
hierarchical coherence ($\rho \approx 1$:1, Fiedler~${\approx}\,0.08$),
and connectivity collapse ($\rho \approx 10^{-5}$:1).

\medskip
\noindent We make the following contributions:

\begin{enumerate}
  \item \textbf{Multi-polytope topology programming.} The Steersman
    programs octahedral ($2{+}2$ at $n{=}4$), RD ($3{+}3$ at $n{=}6$),
    and tesseract ($4{+}4$ at $n{=}8$) geometry into bridge matrices
    using identical hyperparameters. Only the pair specification differs.

  \item \textbf{Spectral attractor universality.} Spectral-only
    training converges to Fiedler ${\approx}\,0.09$ across
    $n \in \{3, 4, 6, 8, 12\}$---a 19.4\% band across a $4\times$
    range of channel counts---establishing the bridge's unprogrammed
    equilibrium.

  \item \textbf{Geometric coherence requirement.} Wrong-labels control
    (random 3-pair partition on $n{=}6$) produces a clean $3{+}3$
    eigenvalue split with co/cross~${\sim}\,10^{-5}$:1, demonstrating
    that eigenvalue structure alone does not constitute topology
    programming. Directional content requires geometrically coherent
    pair specifications.

  \item \textbf{Number-theoretic negative result.} Prime-derived pairs
    (Sophie Germain, twin primes, residue classes) produce chain-like
    eigenvalue progressions, not block structure. Algebraic structure
    without geometric embedding is functionally equivalent to random
    noise ($<10\%$ trajectory divergence from wrong-labels across
    10{,}000 steps).

  \item \textbf{Emanation architecture.} A hierarchical master-plus-offsets
    bridge converges to the spectral attractor rather than block-diagonal
    structure, establishing that topology programming requires direct
    contrastive feedback to individual bridge matrices.

  \item \textbf{Reproducibility.} Tesseract contrastive ($n{=}8$)
    replicates with Pearson $r = 1.0000$ across 34 matching checkpoints
    from independent runs.

  \item \textbf{Four-regime taxonomy.} The complete experimental
    programme yields four dynamical regimes---block-diagonal, spectral
    attractor, wrong-partition, and collapse---cleanly separated with no
    intermediate outcomes.

  \item \textbf{Bridge initialization independence.} Three corpus-coupled
    initializations with maximally different sign patterns converge to
    statistically indistinguishable block-diagonal strength
    (50--52{,}000:1), confirming that topology is determined by the pair
    specification, not the initial conditions.
\end{enumerate}

The remainder of this paper is organized as follows.
\S\ref{sec:background} reviews Paper~3's findings and related work on
programmable inductive biases. \S\ref{sec:method} introduces the
generalized pair specification interface and the experimental design.
\S\ref{sec:three-geometries} presents the multi-polytope programming
results. \S\ref{sec:spectral-attractor} characterizes the spectral
attractor. \S\ref{sec:negative-controls} tests the compiler's rejection
boundary through wrong-labels and resonance controls.
\S\ref{sec:discussion} discusses implications, limitations, and the
four-regime taxonomy. \S\ref{sec:conclusion} concludes.


% ===================================================================
\section{Background}
\label{sec:background}

\subsection{The Learnable Bridge and the Steersman}

Paper~3 in this series~\cite{bielec2026bridge} introduced the \emph{learnable bridge}:
a per-layer coupling matrix $\mathbf{B}_l \in \mathbb{R}^{n \times n}$ that
mediates interaction between $n$ parallel LoRA channels in a multi-channel
adapter. The bridge is trained jointly with the adapter weights, governed by
a cybernetic feedback loop called the \emph{Steersman}.

The Steersman implements a classical control cycle: a \emph{sensor} extracts
the bridge's graph Laplacian eigenvalues at each feedback interval; an
\emph{oscilloscope} tracks connectivity (Fiedler algebraic connectivity
$\lambda_2$) and co-planar/cross-planar coupling ratios; a \emph{steersman}
module adjusts loss weights; and the modified loss signal acts as the
\emph{actuator} on the bridge parameters. The full loop operates at a
configurable interval (default: every 100 optimizer steps) without
interrupting gradient flow.

Two loss components provide independent control channels:
\begin{itemize}
  \item \textbf{Spectral regularization} targets aggregate connectivity by
    penalizing deviations of $\lambda_2$ from a reference value, ensuring the
    bridge develops non-trivial coupling.
  \item \textbf{Contrastive loss} specifies \emph{which} channels should
    couple preferentially: designated \emph{co-axial} pairs are attracted
    (their coupling is reinforced) while \emph{cross-axial} pairs are repelled
    (their coupling is suppressed).
\end{itemize}

Paper~3 demonstrated that this mechanism, applied with pair specifications
derived from the rhombic dodecahedron's 3-axis face-pair geometry at $n{=}6$,
produces block-diagonal bridge structure with co-planar/cross-planar coupling
ratios exceeding 70{,}000:1. The present paper asks whether this result
generalizes beyond a single polytope.

\subsection{Polytope Geometry and Antipodal Pairing}

A \emph{regular polytope} is a geometric object whose vertices, edges, and
faces are congruent under its symmetry group. Several regular and semi-regular
polytopes admit a natural \emph{antipodal pairing}: for each vertex $v$,
there exists a unique vertex $-v$ at maximal distance through the polytope's
center. When the polytope's vertices are mapped to adapter channels, antipodal
pairs define co-axial specifications for the contrastive loss.

We test four polytope geometries spanning 3D and 4D:

\begin{description}
  \item[Octahedron] ($n{=}4$, 2 axes): The regular octahedron in
    $\mathbb{R}^3$ has 6 vertices forming 3 antipodal pairs. We use 4
    vertices (2 pairs) corresponding to a 2-axis partition.
  \item[Rhombic dodecahedron] ($n{=}6$, 3 axes): The RD's 14 vertices
    include 6 degree-4 vertices forming 3 face-diagonal pairs aligned with
    the coordinate axes.
  \item[Tesseract] ($n{=}8$, 4 axes): The 4D hypercube's 16 vertices form
    8 antipodal pairs. We use 8 vertices (4 pairs) corresponding to the four
    coordinate axes.
  \item[24-cell] ($n{=}24$, 12 axes): The $D_4$ root polytope in
    $\mathbb{R}^4$, self-dual, with 24 vertices as permutations of
    $(\pm 1, \pm 1, 0, 0)$. Its 12 antipodal pairs group into 6
    coordinate-plane families (wx, wy, wz, xy, xz, yz).
\end{description}

The key insight is that the Steersman mechanism requires \emph{no
modification} to accommodate different polytopes. Only the pair specification
--- which channels are co-axial --- changes.

\subsection{Programmable Inductive Biases}

Prior work on structured adapters has explored low-rank
variants~\cite{valipour2023dylora,loraxs2024},
mixture-of-experts~\cite{zadouri2023pushing},
and hypernetworks~\cite{ha2017hypernetworks} that generate adapter weights
conditioned on task embeddings. Zhu et al.~\cite{zhu2024asymmetry} analyze
the asymmetric roles of the $A$ and $B$ projections in LoRA, showing that
structure within the adapter matters beyond rank alone.

Recent work has converged independently on block-diagonal adapter structure.
G{\"u}rses et al.~\cite{gurses2026diablo} prove that block-diagonal
fine-tuning (DiaBlo) is theoretically at least as expressive as standard
LoRA and empirically competitive, validating the block-diagonal form as a
first-class adapter architecture. Wang et al.~\cite{wang2025gralora}
(GraLoRA) and Barazandeh et al.~\cite{barazandeh2025localized} (Localized
LoRA) partition adapter matrices into independent blocks for
parameter-efficient fine-tuning. Tian et al.~\cite{tian2026spectral}
(Spectral Surgery) analyze LoRA's spectral properties for training-free
refinement, demonstrating the diagnostic value of spectral decomposition
in adapter weight spaces.

Our approach differs from all of these in a specific way: we do not merely
\emph{impose} block-diagonal structure (as DiaBlo and GraLoRA do through
architectural constraint) but \emph{program} it through cybernetic
feedback. The pair specification declares which channels are co-axial;
the Steersman compiles this specification into the bridge matrix during
training. The topology is a design choice, not an emergent property or
a fixed architectural constraint. This aligns with geometric deep
learning~\cite{bronstein2021geometric}, where the structure of the
computation graph reflects the structure of the problem domain---but
extends it from the network architecture to the adapter's internal
coupling topology.

\subsection{Spectral Graph Theory Preliminaries}

We characterize bridge structure through the graph Laplacian
$\mathbf{L} = \mathbf{D} - \mathbf{W}$, where $\mathbf{W}$ is the
absolute-value weight matrix derived from the bridge and $\mathbf{D}$ is
the degree matrix. The eigenvalues $0 = \lambda_1 \leq \lambda_2 \leq
\cdots \leq \lambda_n$ encode connectivity:

\begin{itemize}
  \item $\lambda_2$ (Fiedler value): algebraic connectivity. $\lambda_2 = 0$
    indicates a disconnected graph; $\lambda_2 > 0$ indicates a connected
    graph.
  \item Block-diagonal structure manifests as $k$ near-zero eigenvalues
    followed by $k$ large eigenvalues, where $k$ equals the number of
    co-axial pairs --- a clean $k{+}k$ split.
  \item The \emph{co-planar/cross-planar ratio} (co/cross) measures
    directional preference: the ratio of mean absolute coupling within
    co-axial pairs to mean absolute coupling between non-paired channels.
\end{itemize}

These two metrics --- Fiedler value and co/cross ratio --- jointly classify
bridge states into the four regimes described in~\S\ref{sec:negative-controls}.

Tam et al.~\cite{tam2020fiedler} introduced differentiable Fiedler
regularization as a sparsity-promoting objective for neural network
pruning, minimizing $\lambda_2$ to disconnect subgraphs. Our spectral
loss inverts this application: we \emph{target} a specific $\lambda_2$
attractor to program topology, not to prune it.


% ===================================================================
% ===================================================================
\section{Method: Generalized Pair Specification}
\label{sec:method}

The Steersman's contrastive loss, as defined in
Paper~3~\cite{bielec2026bridge}, encodes the rhombic dodecahedron's
face-pair geometry through a fixed partition of channel indices into
co-planar and cross-planar sets.  We generalize this mechanism by
replacing the hard-coded RD partition with a configurable pair
specification interface.  The Steersman itself---its three control laws
(connectivity trending, directionality stagnation, stability
monitoring), spectral regularization target, and feedback dynamics---is
\emph{unchanged}.  The only modification is the pair specification fed
to the contrastive loss.

\subsection{The Pair Specification Interface}
\label{sec:pairs}

The function \texttt{\_compute\_pair\_indices}$(n)$ returns two
lists: \emph{co-axial} pairs $\mathcal{C}$ (channels that should
couple strongly) and \emph{cross-axial} pairs $\mathcal{X}$ (all
remaining distinct-index pairs).  The contrastive loss attracts
co-axial pairs and repels cross-axial pairs, operating on the
partition $(\mathcal{C}, \mathcal{X})$ identically to Paper~3.
The spectral loss maintains overall algebraic connectivity by
targeting a Fiedler eigenvalue $\lambda_2 \approx 0.1$.

Four geometric specifications derived from polytope vertex-opposition
structure are tested:

\paragraph{Octahedron ($n{=}4$).}
Two co-axial pairs encoding a 2-axis partition of the octahedron's
three vertex-opposition axes:
\begin{equation}
  \texttt{pairs}_{\text{oct}}(4) = \{(0,1),\; (2,3)\}
  \label{eq:oct_pairs}
\end{equation}
with $|\mathcal{C}| = 2$ and $|\mathcal{X}| = 4$.  This is the
minimum-channel-count geometry: only two co-axial pairs and the
smallest cross/co ratio ($4/2 = 2$), providing the strongest co-axial
signal per channel pair.

\paragraph{Rhombic dodecahedron ($n{=}6$).}
Three co-planar pairs from the RD's 3-axis face-pair geometry:
\begin{equation}
  \texttt{pairs}_{\text{RD}}(6) = \{(0,5),\; (1,4),\; (2,3)\}
  \label{eq:rd_pairs}
\end{equation}
with $|\mathcal{C}| = 3$ and $|\mathcal{X}| = 12$.  Each pair shares
a vanishing coordinate in the face-normal vectors, producing three
independent $2 \times 2$ blocks when compiled.

\paragraph{Tesseract ($n{=}8$).}
Four co-axial pairs aligned with the four coordinate axes of the
4D hypercube:
\begin{equation}
  \texttt{pairs}_{\text{tess}}(8) = \{(0,1),\; (2,3),\; (4,5),\;
    (6,7)\}
  \label{eq:tess_pairs}
\end{equation}
with $|\mathcal{C}| = 4$ and $|\mathcal{X}| = 24$.  If the
Steersman generalizes beyond 3D, this specification should produce
four independent $2 \times 2$ blocks---a $4{+}4$ eigenvalue split.

\paragraph{24-cell ($n{=}24$, D4 root polytope).}
Twelve co-axial pairs from the 24 vertices expressed as permutations
of $(\pm 1, \pm 1, 0, 0)$ in $\mathbb{R}^4$.  Vertices are indexed
by coordinate plane: $wx$ (indices 0--3), $wy$ (4--7), $wz$ (8--11),
$xy$ (12--15), $xz$ (16--19), $yz$ (20--23), with two antipodal
pairs per plane:
\begin{equation}
  \texttt{pairs}_{\text{24c}}(24) =
    \bigl\{(0,3),(1,2),\; (4,7),(5,6),\; (8,11),(9,10),\;
    (12,15),(13,14),\; (16,19),(17,18),\; (20,23),(21,22)\bigr\}
  \label{eq:24c_pairs}
\end{equation}
with $|\mathcal{C}| = 12$ and $|\mathcal{X}| = \binom{24}{2} - 12
= 264$.  The 24-cell is the densest 4D sphere-packing polytope
(the D4 root system), standing in the same relation to 4D as the
rhombic dodecahedron does to 3D.  If the Steersman scales, this
specification should produce twelve independent $2 \times 2$
blocks---a $12{+}12$ eigenvalue split.

The interface imposes no constraint on which indices are paired.
Any partition of $n$~channels into $\lfloor n/2 \rfloor$ pairs
constitutes a valid input.  This generality enables the non-geometric
controls described next.

\subsection{Non-Geometric Controls}
\label{sec:controls}

Two negative controls test whether the Steersman requires geometric
coherence in the pair specification or compiles arbitrary partitions.

\paragraph{Wrong-Labels (WL-001).}
Six channels partitioned into three pairs by random permutation
(seed~42), producing co-axial pairs that do not correspond to any
polytope's antipodal structure.  The Steersman receives the same
contrastive loss function with the same hyperparameters as the
standard $n{=}6$ RD experiment; only the pair specification differs.
If the Steersman is a brute-force topology programmer, random
pairs should produce block-diagonal structure.  If geometric
coherence is required, the random partition should fail.

\paragraph{Resonance (R-001).}
Six channels with co-axial pairs derived from number-theoretic
relationships among corpus primes:  Sophie~Germain
($2(11){+}1 {=} 23 \to$ channels $(0,2)$), consecutive primes
($29, 31 \to$ channels $(3,4)$), and shared residue class
($11 \equiv 17 \equiv 5 \pmod{6} \to$ channels $(0,5)$).  These
pairs carry genuine algebraic structure---prime-theoretic bonds---but
no geometric embedding in any spatial dimension.  This control tests
whether mathematical structure in general, or spatial structure
specifically, is the operative requirement.

\subsection{Emanation Architecture (E-001)}
\label{sec:emanation_method}

The emanation architecture introduces hierarchical bridge structure as
a variant of the standard per-layer independent bridges.  A single
master bridge matrix $\mathbf{M}_{\text{master}}$ is shared across all
transformer layers, with per-layer offset matrices capturing local
variation:
\begin{equation}
  \mathbf{M}^{(l)} = \mathbf{M}_{\text{master}} +
    \boldsymbol{\Delta}^{(l)},
  \label{eq:emanation}
\end{equation}
where offsets are scaled by $1/\sqrt{l}$ for layer index~$l$.  The
Steersman monitors a \emph{coherence} metric---the mean Frobenius norm
of the offsets relative to the master---and applies regularization
pressure when individual layers fragment excessively from the global
topology.  The same RD co-axial pairs are specified, but contrastive
loss operates through the master matrix rather than per-layer bridges.
This tests whether global topological coherence can coexist with local
layer-wise adaptation.

\subsection{Spectral-Only Baseline}
\label{sec:spectral_baseline}

Channel counts $n \in \{3, 4, 6, 8, 12\}$ are trained with spectral
regularization alone---contrastive loss disabled.  The Steersman's
feedback loop targets Fiedler algebraic connectivity
$\lambda_2 \approx 0.1$ but provides no directional preference
among channel pairs.  This baseline establishes the \emph{spectral
attractor}: the bridge's unprogrammed equilibrium, the connectivity
level it reaches when the Steersman provides coupling pressure but
no topological specification.

\subsection{Training Setup}
\label{sec:training_setup}

All experiments use TinyLlama-1.1B-Chat with identical hyperparameters:
learning rate $2 \times 10^{-4}$, batch size~2, gradient
accumulation~8 (effective batch~16), LoRA rank~24, identity
initialization, random seed~42, maximum sequence length~512,
alpaca-cleaned dataset, 10{,}000~training steps.  The channel size
$s = r/n$ varies with~$n$: $s{=}8$ ($n{=}3$), $s{=}6$ ($n{=}4$),
$s{=}4$ ($n{=}6$), $s{=}3$ ($n{=}8$), $s{=}1$ ($n{=}24$).
Table~\ref{tab:topology_specs} summarizes all pair specifications.

\begin{table}[t]
\centering
\caption{Pair specifications for all experimental topologies.  Co-axial
  pairs receive positive contrastive gradient; all remaining
  off-diagonal pairs are cross-axial and receive negative gradient.
  $|\mathcal{C}|$ = co-axial pair count,
  $|\mathcal{X}|$ = cross-axial pair count.}
\label{tab:topology_specs}
\small
\begin{tabular}{llclcc}
\toprule
Topology & $n$ & Co-axial pairs & Source & $|\mathcal{C}|$ & $|\mathcal{X}|$ \\
\midrule
Octahedron (3D) & 4 & $(0,1), (2,3)$ & Vertex opposition & 2 & 4 \\
RD (3D)         & 6 & $(0,5), (1,4), (2,3)$ & Face opposition & 3 & 12 \\
Tesseract (4D)  & 8 & $(0,1), (2,3), (4,5), (6,7)$ & Vertex opposition & 4 & 24 \\
24-cell (4D)    & 24 & 12 antipodal pairs (Eq.~\ref{eq:24c_pairs}) & D4 root system & 12 & 264 \\
\midrule
Wrong-labels    & 6 & Random partition (seed 42) & None (control)  & 3 & 12 \\
Resonance       & 6 & $(0,2), (3,4), (0,5)$ & Number-theoretic & 3 & 12 \\
\bottomrule
\end{tabular}
\end{table}


% ===================================================================
\section{Topology Programming Across Three Geometries}
\label{sec:three-geometries}

The Steersman mechanism introduced in Paper~3 discovered block-diagonal structure
aligned with the rhombic dodecahedron's 3-axis coordinate geometry at $n{=}6$.
We now show this is not specific to the RD\@. The same cybernetic feedback loop,
with no modification except the pair specification, programs topology from three
distinct polytopes: the octahedron ($n{=}4$, 2 axes), the rhombic dodecahedron
($n{=}6$, 3 axes), and the tesseract ($n{=}8$, 4 axes).

\subsection{Experimental Design}

Each experiment uses TinyLlama~1.1B with identical hyperparameters (rank~24,
$\text{lr}{=}2{\times}10^{-4}$, batch~size~2, gradient accumulation~8,
10{,}000 steps). The \emph{only} variable is the contrastive pair specification:

\begin{itemize}
  \item \textbf{Octahedron (O-001, $n{=}4$):} Two co-axial pairs
    $\{(0,1),(2,3)\}$ encoding a 2-axis partition. Four cross-axial pairs.
  \item \textbf{Rhombic dodecahedron (H-ch6, $n{=}6$):} Three co-planar pairs
    from RD face-pair geometry. Twelve cross-planar pairs.
  \item \textbf{Tesseract (T-001r2, $n{=}8$):} Four co-axial pairs
    $\{(0,1),(2,3),(4,5),(6,7)\}$ encoding the four coordinate axes of the
    4D hypercube. Twenty-four cross-axial pairs.
\end{itemize}

\subsection{Results}

All three geometries produce block-diagonal structure with extreme co-planar
vs.\ cross-planar separation:

\begin{table}[h]
\centering
\begin{tabular}{lccccc}
\toprule
\textbf{Geometry} & $n$ & \textbf{Axes} & \textbf{Co/Cross} & \textbf{Fiedler} & \textbf{Eigenvalue Pattern} \\
\midrule
Octahedron  & 4 & 2 & $\sim$10\textsuperscript{5}:1 & $1.1{\times}10^{-5}$ & $[0, 0, 2.06, 2.08]$ \\
RD          & 6 & 3 & 70{,}404:1 & $9.0{\times}10^{-5}$ & $[0, 0, 0, 2.4, 2.4, 2.4]$ \\
Tesseract   & 8 & 4 & 41{,}564:1 & $1.9{\times}10^{-4}$ & $[0, 0, 0, 0, 0.65, 0.68, 0.68, 0.68]$ \\
\bottomrule
\end{tabular}
\caption{Block-diagonal emergence across three polytope geometries. All
experiments use identical hyperparameters; only the contrastive pair
specification differs. Each geometry produces the predicted block structure:
$k{+}k$ eigenvalue split where $k$ equals the number of co-axial pairs.}
\label{tab:three-geometries}
\end{table}

The eigenvalue spectra confirm the predicted structure in each case. The
octahedron produces a clean $2{+}2$ split, the RD a $3{+}3$ split, and the
tesseract a $4{+}4$ split. In all cases, cross-planar eigenvalues are
suppressed by 4--5 orders of magnitude relative to co-planar eigenvalues.

\begin{figure}[t]
  \centering
  \includegraphics[width=\textwidth]{paper4/figures/fig_inverse_n.pdf}
  \caption{Co-planar/cross-planar coupling ratio across polytope geometries.
  The inverse relationship between channel count $n$ and co/cross ratio
  suggests that lower-dimensional polytopes (fewer cross-axial pairs to
  suppress) achieve stronger block separation. Within the RD geometry ($n{=}6$),
  three independent runs produce ratios within 5\% of each other.}
  \label{fig:inverse-n}
\end{figure}

\begin{figure}[t]
  \centering
  \includegraphics[width=\textwidth]{paper4/figures/fig_o001_emergence.pdf}
  \caption{O-001 block-diagonal emergence (octahedron, $n{=}4$). \textbf{Left:}
  Co-planar coupling grows monotonically while cross-planar coupling is
  suppressed by 5 orders of magnitude. \textbf{Right:} The co/cross ratio
  reaches 473{,}622:1 at step 10{,}000, the highest of any geometry tested.}
  \label{fig:o001-emergence}
\end{figure}

\subsection{Interpretation}

The Steersman does not discover geometry---it \emph{compiles} it. Given a
valid pair specification encoding the antipodal structure of a polytope, the
cybernetic feedback loop produces the corresponding block-diagonal partition.
The dimensionality of the source polytope (3D for octahedron and RD, 4D for
tesseract) is irrelevant; what matters is the co-axial pair structure.

This suggests a general principle: \emph{any} polytope whose vertices admit
a natural antipodal pairing can serve as a topology specification for the
Steersman. The bridge matrix is a programmable substrate that accepts
geometric instructions and produces the corresponding algebraic structure.


\subsection{The Tesseract: 4D Programming}
\label{sec:tesseract}

The tesseract (4D hypercube) provides the first test beyond 3D polytopes.
With $n{=}8$ channels and 4 co-axial pairs, the predicted outcome is a
$4{+}4$ eigenvalue split --- four near-zero eigenvalues (one per co-axial
pair) and four large eigenvalues (the cross-axial coupling).

\textbf{Result (T-001r2, 10{,}000 steps):} Co/cross ratio 41{,}564:1.
Fiedler $1.91 \times 10^{-4}$. Eigenvalue pattern:
$[0, 0, 0, 0, 0.65, 0.68, 0.68, 0.68]$ --- a clean $4{+}4$ split.
Validation loss 0.4016, within 0.17\% of baseline. The Steersman programs
4D geometry with the same efficacy as 3D.

\textbf{Reproducibility (T-001r1 vs.\ T-001r2):} Two independent runs with
identical seeds and hyperparameters produce Pearson $r = 1.0000$ at all 6
matching checkpoint steps, with maximum deviation 3.5\%. The cybernetic
training loop is deterministic to within floating-point noise.

\subsection{The 24-Cell: Dense 4D Sphere Packing}
\label{sec:24cell}

The 24-cell is the $D_4$ root polytope --- the unique regular self-dual
polytope in 4D, with 24 vertices as permutations of $(\pm 1, \pm 1, 0, 0)$
in $\mathbb{R}^4$. It tiles 4-space and defines the densest 4D sphere
packing, making it the 4D analogue of the rhombic dodecahedron's role in 3D
(FCC lattice). Its 12 antipodal pairs group into 6 coordinate-plane families
of 2 pairs each, producing $n{=}24$ channels with the richest pair structure
tested.

\textbf{Prediction:} A $12{+}12$ eigenvalue split --- twelve near-zero
eigenvalues (one per co-axial pair) and twelve large eigenvalues. If
confirmed, this demonstrates topology programming at a scale ($24 \times 24$
bridge matrices, 576 learnable parameters per layer) far beyond the $6
\times 6$ bridges of Paper~3.

\textbf{Design note:} At rank 24 with $n{=}24$ channels, each channel
occupies a 1-dimensional subspace of the adapter's rank space. This
extreme partitioning tests whether the Steersman can maintain block-diagonal
coherence when each channel's ``budget'' is minimal. In contrast, T-001r2
($n{=}8$, rank 24) gives each channel a 3-dimensional subspace, and H-ch6
($n{=}6$, rank 24) gives each channel 4 dimensions.

\textbf{Signal density challenge:} The 24-cell's co-axial/cross-axial ratio
is 12/264 $= 0.045$, significantly lower than the tesseract's 4/24 $= 0.167$
(a factor of $3.7\times$). Each contrastive update must reinforce 12 co-axial
pairs while suppressing 264 cross-axial pairs, compared to 4 and 24 for the
tesseract. Whether the contrastive signal remains sufficient at this density
is the central experimental question.

\textbf{Result (24C-001, 10{,}000 steps):} The eigenvalue
spectrum shows a clean $12{+}12$ split emerging from step~300 and growing
monotonically. By step~1{,}000, the low block [0:12] occupies
$[0.000, 0.016]$ while the high block [12:24] clusters at $[0.263, 0.300]$,
producing block separation of $16.8\times$. By step~3{,}000 (30\% of
training), block separation reaches $127\times$ as the high block migrates
to $[0.737, 0.781]$ while the low block compresses to $[0.000, 0.006]$.

Post-hoc co/cross recovery (PC-001) from all 100 saved bridge checkpoints
fills a logging blind spot that had hidden 85\% of the co/cross growth
trajectory. A bug in Control Law~2's $n{=}6$ gate prevented real-time
co/cross tracking for $n{=}24$, but all bridge matrices were saved.
Recovery reveals monotonic growth from $3$:1 (step~100) through
$1{,}832$:1 (step~1{,}000), $5{,}673$:1 (step~4{,}300), $10{,}831$:1
(step~6{,}000), $22{,}455$:1 (step~7{,}500), to a stabilization band of
$34{,}600$--$37{,}600$:1 from step~8{,}000 onward. Final co/cross:
$35{,}808$:1 at step~10{,}000.

The Fiedler trajectory reveals two distinct decline phases. Phase~1
(steps~0--1{,}200): exponential decline from $0.043$ to ${\sim}0.002$,
followed by apparent stabilization near $0.002$ for ${\sim}3{,}000$ steps.
Phase~2 (steps~4{,}000--10{,}000): a second, slower decline carrying the
Fiedler from $0.0026$ through $0.0012$ (step~7{,}000) and $0.00098$
(step~7{,}600, breaking sub-milliunit) to a stabilization band of
$0.000535$--$0.000588$ (steps~8{,}000--10{,}000, confirmed for 2{,}100
steps). The two-phase structure mirrors CW-001's plateau-then-breakout
pattern (\S\ref{sec:disc-whisper}), though here the second decline is gradual
rather than exponential.

At step~10{,}000, the Fiedler has stabilized at
$5.55 \times 10^{-4}$---comparable to CW-001's final value ($4.6 \times
10^{-4}$ at $c_w{=}0.005$) and well above H-ch6's final value ($9 \times
10^{-5}$ at adaptive $c_w$ from $0.1$). The 2{,}100-step stabilization
band confirms the Fiedler floor depends on both $c_w$ and $n$: the fixed
$c_w{=}0.1$ drives strong BD at $n{=}24$ but does not reach the same
depth as adaptive decay at $n{=}6$.

\textbf{Convergence rate scaling.} Comparing co/cross ratios at step~1{,}000
across all four polytope experiments reveals a sharply non-linear
relationship with the per-axis suppression load $\ell = n - 2$ (the number
of cross-axial pairs each co-axial pair must compete against). For $n{=}4$
($\ell{=}2$) and $n{=}6$ ($\ell{=}4$), the step-1{,}000 co/cross ratios are
nearly identical: 7{,}224:1 and 7{,}246:1. At $n{=}8$ ($\ell{=}6$), the
ratio drops 36\% to 4{,}611:1. At $n{=}24$ ($\ell{=}22$), the
step-1{,}000 ratio is $1{,}832$:1 (PC-001 recovery)---a 75\% reduction
from $n{=}6$. However, the 24-cell eventually reaches $35{,}808$:1 by
step~10{,}000, demonstrating that the suppression load slows convergence
\emph{rate} without limiting the eventual \emph{depth}. The ceiling appears
governed by $c_w$, not $\ell$.

\textbf{Confirmation: adaptive vs.\ fixed $c_w$ (24C-002).} A matched
experiment with the Steersman's adaptive decay active (rather than fixed
$c_w{=}0.1$) provides the decisive test. The adaptive mechanism, tuned at
$n{=}6$, decays $c_w$ from $0.1$ to a settled value of ${\sim}0.025$---the
same regime as FI-004's optimal fixed $c_w$ for $n{=}6$. At $n{=}24$,
however, 264 cross-axial pairs require proportionally more contrastive
pressure per co-axial pair than the 4 cross-axial pairs at $n{=}6$. The
adaptive mechanism, calibrated for the latter, \emph{undershoots} for the
former.

The result is unambiguous: final co/cross $5{,}983$:1 (peak $6{,}691$:1
at step~9{,}800)---a factor of $6\times$ below 24C-001's $35{,}808$:1 at
matched step count and identical architecture. The Fiedler eigenvalue
confirms the gap: $0.00461$ (24C-002) vs.\ $0.000555$ (24C-001), an
$8.3\times$ difference. The trajectory is also qualitatively distinct:
24C-002 exhibits a non-monotonic co/cross dip to $936$:1 near step~3{,}000
before recovering---a pattern absent from the fixed-$c_w$ run, likely
reflecting the Steersman's overshoot-then-correct dynamics at low $c_w$.

This confirms the ceiling claim: at $n{=}24$, reducing $c_w$ from $0.1$
to $0.025$ reduces the final BD strength by $6\times$. The contrastive
weight governs the \emph{depth} of block-diagonal formation, not merely
its \emph{speed}---resolving the speed-vs.-attractor question posed in
\S\ref{sec:disc-whisper} decisively in favour of the attractor hypothesis.

\begin{figure}[t]
  \centering
  \includegraphics[width=\textwidth]{paper4/figures/fig_24cell_emergence.pdf}
  \caption{24-cell experiment (24C-001) full trajectory, 10{,}000 steps.
  (a)~Fiedler value (log scale) showing two-phase decline: exponential to
  ${\sim}0.002$ (steps 0--1{,}200), then gradual to stabilization band
  $0.000535$--$0.000588$ (steps 8{,}000--10{,}000, 2{,}100 steps).
  (b)~Co/cross ratio (PC-001 post-hoc recovery) reaching $35{,}808$:1
  at step~10{,}000. Block separation $127\times$ at 30\% of training.
  Fixed $c_w{=}0.1$.}
  \label{fig:24cell-emergence}
\end{figure}

% NOTE: 24C-001 COMPLETE at step 10000. PC-001 recovery DONE — co/cross 35,808:1 final.
% CL2 logging bug: co_cross gated to n==6 in train_contrastive_bridge.py logging code.
% Contrastive LOSS was applied correctly via train_cybernetic.py (no n==6 gate on loss).
% PC-001 results saved to 24C-001/pc_001_co_cross_recovery.json (101 checkpoints).


% ===================================================================
\section{The Spectral Attractor}
\label{sec:spectral-attractor}

\subsection{Universal Convergence}

When the Steersman operates with spectral regularization alone---no contrastive
loss specifying which channels should couple preferentially---the bridge
converges to a characteristic connectivity level independent of the number of
channels. We term this the \emph{spectral attractor}.

\begin{table}[h]
\centering
\begin{tabular}{lccc}
\toprule
\textbf{Experiment} & $n$ & \textbf{Final Fiedler} & \textbf{Topology} \\
\midrule
H-ch3  & 3  & 0.0951 & spectral-only \\
H-ch4  & 4  & 0.0918 & spectral-only \\
E-001  & 6  & 0.0836 & emanation (spectral effective) \\
H-ch8  & 8  & 0.0944 & spectral-only \\
H-ch12 & 12 & 0.1019 & spectral-only \\
\midrule
\multicolumn{2}{l}{\textbf{Mean $\pm$ SD}} & $0.0934 \pm 0.0065$ & \\
\multicolumn{2}{l}{\textbf{Band}} & 19.4\% (min/max) & \\
\bottomrule
\end{tabular}
\caption{Spectral attractor convergence across channel counts $n \in \{3, 4, 6, 8, 12\}$.
All experiments converge to Fiedler $\approx 0.09$ regardless of $n$.
E-001 (emanation architecture) converges to the same attractor, confirming
that hierarchical constraints do not override the spectral equilibrium.}
\label{tab:spectral-attractor}
\end{table}

The spectral attractor at Fiedler $\approx 0.09$ is stable across a $4\times$
range of channel counts ($n{=}3$ to $n{=}12$), with only 19.4\% variation
(Figure~\ref{fig:spectral-attractor}).
This invariance suggests the attractor is a property of the spectral
regularization objective itself, not of the bridge dimensionality.

\begin{figure}[t]
  \centering
  \includegraphics[width=0.7\textwidth]{paper4/figures/fig_spectral_attractor.pdf}
  \caption{Spectral attractor: universal Fiedler convergence across channel
  counts $n \in \{3, 4, 8, 12\}$. The shaded band shows 19.4\% variation
  around the mean (0.0934). The spectral attractor is the bridge's
  unprogrammed connectivity equilibrium.}
  \label{fig:spectral-attractor}
\end{figure}

\subsection{Bifurcation Under Contrastive Loss}

Adding the contrastive loss breaks the spectral attractor. At $n{=}6$ with
RD contrastive pairs, the Fiedler drops from $\sim$0.09 (attractor) to
$9.0 \times 10^{-5}$---a factor of 1{,}000$\times$. The contrastive loss
does not merely redirect connectivity; it \emph{concentrates} it along the
specified axes while suppressing it everywhere else.

This bifurcation reveals the two independent control channels of the
Steersman:
\begin{enumerate}
  \item \textbf{Spectral regularization} provides \emph{connectivity}---how
    much total coupling the bridge develops.
  \item \textbf{Contrastive loss} provides \emph{topology}---which channels
    receive the coupling budget.
\end{enumerate}

Without contrastive specification, connectivity distributes uniformly across
all channel pairs, producing the attractor. With specification, the same
connectivity budget concentrates along the specified axes, producing
block-diagonal structure.

\subsection{Emanation as Spectral Attractor Variant}

E-001 uses the emanation architecture (a single master bridge with per-layer
offsets and coherence monitoring) instead of independent per-layer bridges.
Its final state---Fiedler 0.084, co/cross 1.12---places it squarely in the
spectral attractor regime. The hierarchical constraint maintains layer
coherence (low deviation from master) but does not create directional
preference. Emanation is thus a \emph{constrained} variant of spectral-only
training, not a topology-programming mechanism.


% ===================================================================
\section{Topology Programmability: Negative Controls}
\label{sec:negative-controls}

The three-geometry results (\S\ref{sec:three-geometries}) show the Steersman
compiles valid geometric programs. Two negative controls test the boundaries
of this claim: what happens when the pair specification lacks geometric
coherence?

\subsection{Wrong-Labels (WL-001)}

WL-001 uses a random partition of 6 channels into 3 pairs (seed 42),
producing pairs that do not correspond to any polytope's antipodal structure.
The Steersman receives the same contrastive loss function with the same
hyperparameters---only the pair specification differs.

\textbf{Result:} The eigenvalue spectrum at step 10{,}000 shows a $3{+}3$
split: $[0, 10^{-5}, 2{\times}10^{-5}, 2.437, 2.438, 2.444]$. The
contrastive loss \emph{does} create a 2-partition---any specification of 3
pairs on 6 channels will bisect the channel space. However, the co/cross
ratio is $8.7 \times 10^{-6}$, effectively zero. Compare to H-ch6 (RD
contrastive) at $70{,}404{:}1$.

The interpretation is precise: the contrastive mechanism creates a partition
from \emph{any} 3-pair specification, but only geometrically coherent pairs
produce directional structure within that partition. Wrong-labels produces
two blocks with no internal organization; correct labels produce two blocks
whose internal coupling is $10^5\times$ stronger than their cross-block
coupling.

\subsection{Resonance (R-001)}

R-001 uses pairs derived from number-theoretic relationships among the corpus
primes: Sophie Germain ($2(11){+}1{=}23$, channels 0--2), consecutive primes
($(29,31)$, channels 3--4), and residue class ($11 \equiv 17 \pmod{6}$,
channels 0--5). These pairs carry genuine algebraic structure but no
geometric embedding.

\textbf{Result:} Fiedler $1.2 \times 10^{-5}$, co/cross $\approx 0$.
The eigenvalue pattern is $[0, 2{\times}10^{-5}, 1.222, 2.447, 3.667]$---a
chain-like progression, \emph{not} a clean block split. The prime-theoretic
relationships create a connectivity \emph{gradient} rather than a partition.

This is a stronger negative than WL-001: wrong-labels at least produces the
right \emph{shape} (two blocks) with the wrong content; resonance produces
the wrong shape entirely. Number-theoretic structure is insufficient for
topology programming. The Steersman requires geometric coherence---a valid
antipodal pairing from a polytope's vertex set.

\begin{figure}[t]
  \centering
  \includegraphics[width=\textwidth]{paper4/figures/fig_wl_r001_identity.pdf}
  \caption{Collapse trajectories for the two negative controls. WL-001
  (random pairs) and R-001 (prime-theoretic pairs) produce indistinguishable
  co/cross and Fiedler trajectories, both converging to $\sim$0. The
  mechanisms differ---WL-001 creates blocks without direction; R-001 creates
  a gradient without blocks---but the endpoint is the same: no topology.}
  \label{fig:wl-r001}
\end{figure}

\subsection{Regime Classification}

The complete experimental landscape now partitions into four regimes:

\begin{table}[h]
\centering
\small
\begin{tabular}{llccp{5cm}}
\toprule
\textbf{Regime} & \textbf{Experiments} & \textbf{Fiedler} & \textbf{Co/Cross} & \textbf{Mechanism} \\
\midrule
Block-diagonal & H-ch6, O-001, T-001r2 & $\sim$10\textsuperscript{-5} & $10^4$--$10^5$:1 & Geometric contrastive \\
Spectral attractor & H-ch3/4/8/12, E-001 & $\sim$0.09 & $\sim$1:1 & Spectral-only or non-directional \\
Wrong-partition & WL-001 & $\sim$10\textsuperscript{-5} & $\sim$0 & Contrastive with incoherent pairs \\
Collapse & R-001 & $\sim$10\textsuperscript{-5} & $\sim$0 & Non-geometric pair structure \\
\bottomrule
\end{tabular}
\caption{Four training regimes. Block-diagonal requires geometric contrastive
pairs. The spectral attractor is the bridge's unprogrammed equilibrium.
Wrong-partition creates blocks without directional structure. Collapse
produces neither blocks nor connectivity.}
\label{tab:four-regimes}
\end{table}

The four-regime classification establishes a clear hierarchy: topology
programming requires (1)~the Steersman (cybernetic feedback), (2)~contrastive
loss, and (3)~geometrically coherent pair specification. Remove any one of
these three and the bridge converges to one of the non-BD regimes.


% ===================================================================
\section{The Four Regimes}
\label{sec:four-regimes}

The complete experimental programme --- three polytope geometries, five
spectral-only baselines, two non-geometric controls, and an emanation
architecture --- reveals four distinct training outcomes. Together they form
a taxonomy of bridge behavior under the Steersman.

\subsection{Regime Classification}

\begin{table*}[t]
\centering
\small
\begin{tabular}{llcccp{4.5cm}}
\toprule
\textbf{Regime} & \textbf{Experiments} & \textbf{Fiedler} & \textbf{Co/Cross} & \textbf{Eigenvalues} & \textbf{Mechanism} \\
\midrule
\textbf{1: Block-Diagonal} & H-ch6, O-001, T-001r2 & $\sim$10\textsuperscript{-5}--10\textsuperscript{-4} & 10\textsuperscript{4}--10\textsuperscript{5}:1 & Clean $k{+}k$ split & Geometric contrastive: valid polytope pair spec \\
\textbf{2: Spectral Attractor} & H-ch3/4/8/12, E-001 & $0.084$--$0.102$ & $\sim$1:1 & Smooth distribution & Spectral-only or non-directional architecture \\
\textbf{3: Wrong-Partition} & WL-001 & $\sim$10\textsuperscript{-5} & $\sim$0 & $k{+}k$ split, wrong direction & Contrastive with incoherent pair spec \\
\textbf{4: Collapse} & R-001 & $\sim$10\textsuperscript{-5} & $\sim$0 & Chain-like progression & Non-geometric pair structure \\
\bottomrule
\end{tabular}
\caption{The four training regimes. Block-diagonal (Regime~1) requires all
three ingredients: cybernetic feedback, contrastive loss, and geometrically
coherent pair specification. The spectral attractor (Regime~2) is the
bridge's unprogrammed equilibrium. Wrong-partition (Regime~3) produces
eigenvalue structure without directional content. Collapse (Regime~4) produces
neither blocks nor connectivity.}
\label{tab:four-regimes-full}
\end{table*}

The regimes are distinguished by two independent diagnostics:

\begin{enumerate}
  \item \textbf{Eigenvalue structure:} Does the Laplacian spectrum exhibit a
    clean $k{+}k$ split (near-zero eigenvalues followed by large eigenvalues)?
    Regimes 1 and 3 exhibit this structure; Regimes 2 and 4 do not.
  \item \textbf{Directional coherence:} Does the co-planar/cross-planar
    coupling ratio indicate systematic directional preference? Only Regime~1
    produces this; all other regimes have co/cross $\sim$1 or $\sim$0.
\end{enumerate}

This creates a $2 \times 2$ classification:

\begin{table}[h]
\centering
\begin{tabular}{l|cc}
 & \textbf{Has eigenvalue structure} & \textbf{No structure} \\
\hline
\textbf{Directional} & Regime 1 (BD) & --- \\
\textbf{Non-directional} & Regime 3 (wrong-partition) & Regime 4 (collapse) \\
\end{tabular}
\caption{The $2 \times 2$ regime matrix. Regime~2 (spectral attractor)
precedes this classification --- it is the default state before any pair
specification is applied.}
\label{tab:regime-matrix}
\end{table}

Figure~\ref{fig:four-regimes} illustrates the four regimes with representative
trajectories.

\begin{figure*}[t]
  \centering
  \includegraphics[width=\textwidth]{paper4/figures/fig_four_regimes.pdf}
  \caption{The four regimes of topology programming. \textbf{(a)}~Block-diagonal:
  O-001 co/cross ratio climbs monotonically to 473{,}622:1.
  \textbf{(b)}~Spectral attractor: Fiedler converges to $\sim$0.09 regardless
  of channel count. \textbf{(c)}~Hierarchical coherence (emanation): Fiedler
  and co/cross both converge to the spectral attractor, confirming hierarchy
  provides connectivity without directionality. \textbf{(d)}~Connectivity
  collapse: WL-001 and R-001 produce identical trajectories approaching zero.}
  \label{fig:four-regimes}
\end{figure*}

The fourth cell (directional without eigenvalue structure) is unoccupied in
our experiments. This absence may reflect a geometric constraint: directional
coherence requires the channel space to partition into coupled blocks, which
necessarily produces eigenvalue structure. Whether this cell is geometrically
impossible or merely unvisited remains an open question.

\subsection{Regime Transitions}

The regimes form a hierarchy of requirements. Topology programming
(Regime~1) requires three simultaneous conditions:

\begin{enumerate}
  \item \textbf{Cybernetic feedback} (the Steersman): Without the
    sensor--oscilloscope--steersman--actuator loop, bridge matrices show no
    systematic structure development.
  \item \textbf{Contrastive loss}: Without pair-based attraction/repulsion,
    the bridge converges to the spectral attractor (Regime~2) regardless of
    channel count or architecture.
  \item \textbf{Geometrically coherent pair specification}: Without valid
    polytope geometry, contrastive loss produces either structure without
    direction (Regime~3) or complete collapse (Regime~4).
\end{enumerate}

Removing any single condition produces a predictable degradation:
\begin{itemize}
  \item Remove contrastive loss $\rightarrow$ Regime~2 (spectral attractor)
  \item Replace geometric pairs with random pairs $\rightarrow$ Regime~3
    (wrong-partition)
  \item Replace geometric pairs with algebraic pairs $\rightarrow$ Regime~4
    (collapse)
\end{itemize}

The distinction between Regimes~3 and~4 is diagnostic. Wrong-labels (random
partition) creates valid block structure because \emph{any} partition of $n$
channels into $k$ pairs induces a $k$-way bisection --- the contrastive loss
acts on the partition structure. But the resulting blocks have no internal
orientation: co-axial coupling equals cross-axial coupling within the blocks.
Resonance (number-theoretic) fails even to create clean blocks: the overlapping
pair structure (some channels appear in multiple relationships) produces a
connectivity \emph{gradient} rather than a partition.

\subsection{The Spectral Attractor as Ground State}

Regime~2 occupies a distinguished role: it is the bridge's \emph{ground
state}. Five independent experiments across $n \in \{3, 4, 6, 8, 12\}$ and
one architectural variant (emanation) converge to Fiedler $\approx 0.09$
(\S\ref{sec:spectral-attractor}). The contrastive loss does not add
connectivity to this ground state --- it \emph{redirects} it. The total
coupling budget remains approximately constant; the contrastive loss
concentrates it along the specified axes while suppressing it everywhere
else. This concentration produces the 1{,}130$\times$ Fiedler drop (from
0.09 to $9.0 \times 10^{-5}$) observed when comparing spectral-only to RD
contrastive at $n{=}6$.


% ===================================================================
\section{Discussion}
\label{sec:discussion}

\subsection{The Bridge as Programmable Substrate}
\label{sec:disc-substrate}

Paper~3 framed the bridge as a geometry-\emph{discovering} mechanism:
cybernetic feedback reveals latent rhombic dodecahedral structure in
$6$-channel LoRA bridges. The present paper reframes this. The bridge
is a programmable substrate. The Steersman is the programmer.

The evidence for this reframing is straightforward. Three polytope
geometries---octahedron ($n{=}4$, $2{+}2$ blocks), rhombic dodecahedron
($n{=}6$, $3{+}3$ blocks), and tesseract ($n{=}8$, $4{+}4$
blocks)---all produce block-diagonal structure with coupling ratios
exceeding 40{,}000:1 when the pair specification encodes their
respective co-axial decompositions. The Steersman's control laws,
loss formulation, and hyperparameters remain identical across all three;
only the pair specification $\mathcal{C}$ changes. The bridge compiles
whatever geometric program it receives.

But the compiler is selective. The wrong-labels experiment (WL-001) is
the decisive evidence. A random 3-pair partition of 6~channels produces
a clean $3{+}3$ eigenvalue split---proof that the contrastive mechanism
\emph{can} bisect the channel space from any pair specification. Yet the
co/cross ratio is $8.7 \times 10^{-6}$:1, functionally zero. The
Steersman creates eigenvalue structure from any pair specification, but
it cannot impose directional coherence without genuine geometry. The
distinction between ``structured'' (eigenvalue split) and ``directed''
(co/cross $\gg 1$) is the central finding of this paper, and it
sharpens the characterization of the Steersman from general-purpose
enforcer to selective geometric compiler.

\subsection{Bridge Initialization Independence}
\label{sec:disc-init-independence}

A natural question follows: does the bridge's final topology depend on the
specific initialization, or does the Steersman converge to a canonical
structure regardless of starting conditions?

The FI-002 experiment addresses this directly. Three corpus-coupled
initializations---identical in hexagram structure and geometric coupling
but differing in channel permutation---produce sign patterns with Hamming
distances up to 9/12 (75\% of cross-planar signs differ at initialization).
Despite these differences, all three converge to \emph{identical}
block-diagonal structure: co/cross ratios of 50{,}344:1,
51{,}677:1, and 50{,}382:1 respectively, with validation losses within
0.02\% of each other. The mean bridge matrices agree to four decimal
places (Frobenius distance ${\sim}10^{-4}$), and all 12 consensus sign
positions match perfectly across configurations. Critically, the
configuration with 75\% wrong initial signs (P-001, 9/12 Hamming
distance) converges to the same sign pattern as the configuration
with 100\% correct initial signs (P-000)---the Steersman actively
corrects misaligned signs during training.

A control run with identity initialization (no
corpus coupling) reaches 8{,}085:1 at step~1{,}000---remarkably close to the
original H-ch6 experiment's 7{,}246:1 at the same step, despite the corpus-coupled
run starting from a geometrically structured bridge while the control starts from
identity matrices. By step~1{,}300, P-CTRL reaches 10{,}654:1, but the growth rate
collapses from $+1{,}329$/step to $+98$/step (93\% deceleration) while the Fiedler
value bounces from its minimum ($9.9 \times 10^{-5}$) to $1.07 \times 10^{-4}$---the
same plateau-onset dynamics observed in H-ch6 at comparable training fractions.
The identity control produces random sign patterns (50\% positive, 41.7\%
consensus agreement with corpus configurations), confirming that the
canonical sign pattern requires corpus-coupled initialization to
establish but is independent of the specific channel mapping used.
The early-phase convergence rate is pair-specification-determined.

The convergence of maximally different initial sign patterns to equivalent
final topology has two implications. First, the Steersman is robust to
initialization: the same geometric program compiles to the same
block-diagonal structure regardless of the bridge's starting state. The
topology is determined by the pair specification, not the initial
conditions. Second, the sign pattern frozen into the bridge during
training---which FI-001 established as 98.2\% stable after $\sim$1{,}200
steps---may represent a degree of freedom that is determined by
initialization but irrelevant to the block-diagonal structure itself.
The topology has a well-defined ``shape'' (block-diagonal with $k{+}k$
eigenvalue split) and an initialization-dependent ``orientation'' (sign
pattern within the blocks) that the Steersman does not constrain.

This decomposition---universal topology, initialization-dependent
orientation---parallels the distinction between a crystal's lattice
structure and its orientation in space. The Steersman programs the
lattice; it does not select the orientation.

\begin{figure}[t]
  \centering
  \includegraphics[width=\textwidth]{paper4/figures/fig_fi002_init_independence.pdf}
  \caption{Bridge initialization independence (FI-002). Three corpus-coupled
  initializations with Hamming distances up to 9/12 converge to
  statistically indistinguishable co/cross ratios (50--52{,}000:1).
  An identity-initialized control (P-CTRL) reaches 10{,}654:1 at step~1{,}300,
  following the same $3{+}3$ block-diagonal trajectory with plateau onset
  matching H-ch6's two-phase dynamics.
  Topology is pair-specification-determined, not initialization-determined.}
  \label{fig:fi002-init}
\end{figure}

\subsection{Topology Requires Continuous Maintenance}
\label{sec:disc-homeostasis}

A natural follow-up: once the Steersman has programmed the bridge, is the
topology self-sustaining? The crystal analogy from \S\ref{sec:disc-init-independence}
would predict yes---crystals maintain their lattice without external forcing.
FI-003 tests this directly.

Starting from a fully converged $n{=}6$ bridge (P-000 adapter, co/cross
50{,}344:1 after 3{,}000 Steersman-guided steps), FI-003 continues training
with \emph{only} the language modelling loss---no spectral loss, no
contrastive loss. The Steersman is switched off. If the block-diagonal
topology is an attractor of the LM loss landscape (i.e., the structure is
self-sustaining once established), the topology should persist. If it
requires active maintenance, it should decay.

The topology decays immediately. Within 100~steps of LM-only training (3.3\%
of the original training budget), the sign pattern collapses from 100\%
agreement with the converged snapshot to 50.85\%---statistically
indistinguishable from random (50\%). The co/cross ratio undergoes
exponential decay: $12{,}586 \to 179 \to 60 \to 30 \to 22 \to 18 \to 16
\to 14 \to 13 \to 12 \to 11 \to 10$ over 1{,}200~steps (40\% of the original
training budget), approaching the isotropic equilibrium. The decay
exhibits two timescales: a fast component (half-life ${\sim}15$~steps)
that destroys the block-diagonal structure, and a slow component
(half-life ${\sim}230$~steps) that diffuses the remaining directional
coherence toward isotropy. The bridge magnitudes change minimally
($0.0277 \to 0.0280$), confirming that the decay is not driven by weight
magnitude changes but by loss of directional coherence. Validation loss
rises from 0.499 to 0.608, then partially recovers to
${\sim}0.52$---the adapter still functions as a LoRA module but loses
its geometric organization.

The crystal analogy fails. A better analogy is a spinning top: the
block-diagonal structure is dynamically stable under the Steersman's
continuous forcing but immediately topples to the isotropic equilibrium
when the forcing is removed. The Steersman does not program the bridge
once and walk away; it performs continuous homeostatic maintenance.

\begin{figure}[t]
  \centering
  \includegraphics[width=\textwidth]{paper4/figures/fig_fi003_homeostasis.pdf}
  \caption{Topology homeostasis (FI-003). Starting from a converged
  bridge (co/cross 50{,}344:1), the Steersman is removed and training
  continues with LM loss only. (a)~Co/cross ratio decays exponentially
  from 12{,}586:1 toward the isotropic equilibrium, with decelerating
  halving time. (b)~Sign stability collapses from 100\% to random
  (${\sim}50\%$) within 100~steps---the block-diagonal structure is
  dynamically maintained, not permanently impressed.}
  \label{fig:fi003-homeostasis}
\end{figure}

Combined with \S\ref{sec:disc-init-independence}, this yields a complete
characterization. The Steersman is both robust (same topology regardless
of initialization) and necessary (topology dissolves without it). The
block-diagonal structure exists in a region of parameter space that is
\emph{accessible} from any initialization via the Steersman's loss
landscape, but \emph{unstable} under the LM loss alone. The spectral
attractor at Fiedler~$\approx 0.09$ is the true minimum of the LM-only
landscape; the block-diagonal configuration is a Steersman-maintained
fixed point that requires active control to sustain.

\subsection{Steersman Annealing: Five Regimes of Control}
\label{sec:disc-annealing}

FI-003 established that removing the Steersman entirely causes rapid
topology decay. A natural question follows: how does the topology respond
to \emph{gradual} removal? FI-004 answers this by linearly annealing both
contrastive ($0.1 \to 0$) and spectral ($0.05 \to 0$) weights over
3{,}000~steps from a converged P-000 adapter.

The result is not a monotonic decay. Five distinct regimes emerge as the
contrastive weight $c_w$ decreases:

\begin{enumerate}
\item \textbf{Maintenance} ($c_w = 0.10$): The starting point. Co/cross
  12{,}586:1, inherited from the converged P-000 adapter. Signs 100\%
  stable.
\item \textbf{Interference} ($0.087 < c_w < 0.097$): Rapid co/cross
  collapse to a floor of 7.3:1 within 400~steps. The contrastive and
  LM gradients are in destructive conflict---full-strength Steersman
  forcing fights the optimizer's preferred direction, producing worse
  topology than either alone.
\item \textbf{Growth} ($0.017 < c_w < 0.087$): Exponential co/cross
  recovery from 7.3:1 to a peak of 18{,}671:1 at $c_w = 0.017$
  (step~2{,}500). The growth oscillates between compression
  ($1.17$--$1.22\times$/step) and expansion ($1.37$--$1.55\times$/step)
  phases, with a geometric mean of $\sim$$1.33\times$/100~steps. During this
  regime, sign stability settles to $51$--$55\%$ (random). The weak
  contrastive signal provides a directional bias that the LM optimizer
  amplifies without gradient conflict.
\item \textbf{Decay} ($0.003 < c_w < 0.017$): Gradual decline from
  18{,}671:1 to 12{,}467:1 ($\sim$6\%/step). The contrastive signal
  weakens below the threshold required to maintain the growth trajectory,
  but topology degrades slowly.
\item \textbf{Cliff} ($c_w = 0.003 \to 0.000$): Co/cross drops 76\%
  from 12{,}467:1 to 2{,}942:1 in the final step. The last 0.3\% of
  contrastive weight was maintaining three-quarters of the directional
  coherence.
\end{enumerate}

Two findings are striking. First, the \textbf{optimal operating point
is at $c_w \approx 0.02$}, not at the nominal training weight of $0.10$.
Full-strength Steersman forcing operates in the interference regime.
The topology programmer works best when it whispers. Second, even at
$c_w = 0$, the bridge retains 2{,}942:1---$294\times$ higher than
FI-003's comparable training fraction ($\sim$10:1 at step~1{,}200 of
LM-only training). The annealing process imprints topology that persists
beyond the Steersman's removal, though presumably the residual will decay
under extended LM-only training following FI-003's exponential pattern.

Throughout all five regimes, sign stability remains at $51$--$55\%$
(random). The co/cross ratio reaches 18{,}671:1 with signs statistically
indistinguishable from chance. This confirms the decomposition from
\S\ref{sec:disc-init-independence}: directional coherence (magnitude
alignment across co-axial pairs) and sign structure are orthogonal
properties of the bridge. The annealing experiment programs the former
without the latter.

Validation loss improves monotonically from 0.499 to 0.315 throughout
training, confirming that the topology programming does not degrade
language modelling performance. The Fiedler value rises from 0.10 to
0.47 as the spectral constraint anneals away---the bridge becomes
over-connected as the connectivity regularization disappears, but this
excess connectivity does not impair downstream performance.

\begin{figure}[t]
  \centering
  \includegraphics[width=\textwidth]{paper4/figures/fig_fi004_annealing.pdf}
  \caption{Steersman annealing (FI-004). (a)~Co/cross ratio vs.\ training
  step showing five regimes: maintenance (step~0), interference (floor 7:1),
  exponential growth (peak 18{,}671:1 at $c_w{=}0.017$), gradual decay,
  and cliff ($76\%$ drop at $c_w{=}0$). (b)~Operating curve: co/cross vs.\
  contrastive weight reveals the optimal operating point at
  $c_w \approx 0.02$, not at the nominal training weight of $0.10$.}
  \label{fig:fi004-annealing}
\end{figure}

\subsection{Implications for Adapter Design}
\label{sec:disc-adapter}

If the bridge is programmable, topology becomes a design parameter.

\paragraph{Topology as hyperparameter for PEFT.}
Current parameter-efficient fine-tuning methods treat the adapter's
internal structure as homogeneous or leave it to the optimizer.
DiaBlo~\cite{gurses2026diablo} proves that block-diagonal structure is at
least as expressive as standard LoRA, and GraLoRA~\cite{wang2025gralora}
demonstrates competitive performance with granular block partitions.
These results validate block-diagonal adapters as a structural class.
The Steersman adds a further axis: the \emph{topology specification}---a
set of co-axial pairs derived from a polytope's vertex decomposition.
Where DiaBlo and GraLoRA fix the block structure architecturally, the
Steersman \emph{programs} it through cybernetic feedback, making the
block topology a trainable design parameter rather than a static
constraint. This specification joins rank $r$, learning rate, and target
modules as a hyperparameter that the practitioner selects before training.

\paragraph{Task-specific topology.}
Different task families may benefit from different geometric priors.
A task requiring two independent subspaces (e.g., separating style from
content) might use octahedral topology ($n{=}4$, 2~axes). A task with
higher-dimensional factored structure might use tesseract topology
($n{=}8$, 4~axes). The inverse-$n$ relationship between channel count
and coupling strength (474K:1 at $n{=}4$ vs.\ 42K:1 at $n{=}8$)
provides a practical trade-off: fewer axes yield stronger per-axis
signal, while more axes provide finer-grained topological control.

\paragraph{The four regimes as diagnostic framework.}
The four-regime taxonomy provides a diagnostic for any multi-channel
bridge configuration. Practitioners should measure both the eigenvalue
pattern (to detect whether a partition exists) and the co/cross ratio
$\rho$ (to detect whether the partition carries directional content).
A configuration with a clean eigenvalue split but $\rho \approx 0$
occupies Regime~3 (wrong-partition)---structured but undirected. Only
configurations with both a clean split \emph{and} $\rho \gg 1$ occupy
Regime~1 (block-diagonal). This two-metric diagnostic prevents the
false positive of interpreting eigenvalue structure as evidence of
geometric programming when no directional coherence exists.

\paragraph{Multi-modal adapters.}
When adapting models across modalities, different modalities could carry
different geometric structures within a shared bridge. The pair
specification provides an interface for encoding known cross-modal
relationships---e.g., specifying which channels should couple across
a vision-language boundary---while the Steersman enforces the
specification during training. The block structure would then serve as
both an inductive bias and a compositional interface: blocks
corresponding to different modalities could be independently updated,
merged, or swapped.

\subsection{The Spectral Attractor as Default State}
\label{sec:disc-attractor}

Without geometric programming, bridges reach a universal connectivity
equilibrium: Fiedler $\approx 0.09$, stable across a $4\times$ range
of channel counts ($n \in \{3, 4, 8, 12\}$, 10.4\% band). This is the
bridge's unprogrammed state---connected but undirected, with coupling
distributed isotropically across all channel pairs.

The spectral attractor clarifies what contrastive loss does. It does not
\emph{add} connectivity to the bridge. The bridge already has
connectivity---approximately $\lambda_2 \approx 0.09$ worth of it.
Contrastive loss \emph{redirects} this connectivity from the isotropic
attractor to specific axes. The $1{,}130\times$ Fiedler drop at $n{=}6$
(from $0.09$ to $9.0 \times 10^{-5}$) quantifies this redirection: the
connectivity budget that was distributed across all 15~channel pairs is
concentrated into 3~co-axial pairs, suppressing the remaining 12 by
a factor of $\sim$70{,}000.

The emanation architecture (E-001) reaches the same attractor through a
different mechanism---a master bridge with per-layer offsets and
coherence monitoring. Its convergence to Fiedler $0.084$, co/cross
$1.12$:1, confirms that the spectral attractor is robust to
architectural variation. Whether the bridge is independent per-layer
(spectral-only experiments) or hierarchically constrained (emanation),
the unprogrammed equilibrium is the same. The attractor is a property
of the spectral loss landscape, not of the bridge architecture.

\subsection{What Constitutes Geometric Coherence?}
\label{sec:disc-coherence}

The paper demonstrates that geometric coherence in the pair
specification is required for topology programming, but does not
provide a formal characterization of what makes a specification
``geometric.'' This is an open question.

What has been tested: three confirmed polytope pair specifications
work---octahedron, rhombic dodecahedron, tesseract---and a fourth
(24-cell) shows the predicted $12{+}12$ eigenvalue split with
$127\times$ block separation at 30\% of training and co/cross
$35{,}808$:1 at completion. These share a common property: each encodes the antipodal
vertex structure of a polytope embeddable in Euclidean space, with the
co-axial pairs aligned to coordinate axes. Random partitions that do not
correspond to any polytope's antipodal structure produce eigenvalue
splits without directional content (WL-001). Number-theoretic bonds that
carry genuine algebraic structure but no spatial embedding produce
neither eigenvalue blocks nor directional content (R-001).

A working hypothesis emerges: the pair specification must correspond to
the antipodal structure of a regular or semi-regular polytope
embedded in the channel space, such that the co-axial pairs identify
axis-aligned subspaces that can be independently optimized. This
would explain why four polytopes spanning two spatial dimensionalities
(3D octahedron and RD, 4D tesseract and 24-cell) all succeed---they
share the antipodal property---while random and algebraic specifications
fail.

Several directions remain unexplored. Do irregular polytopes work?
Do non-antipodal pair structures (e.g., face-diagonal rather than
body-diagonal) produce block-diagonal structure? How does the ratio
of co-axial to cross-axial pairs ($|\mathcal{C}|/|\mathcal{X}|$)
interact with the coupling ratio $\rho$?

\subsection{Contrastive Weight as Speed Control}
\label{sec:disc-whisper}

A natural question arises from the four-regime taxonomy: does the
contrastive weight $c_w$ control \emph{whether} block-diagonal structure
emerges, or \emph{how quickly} it emerges? The FI-004 annealing experiment
(\S\ref{sec:disc-annealing}) found an optimal fixed
$c_w \approx 0.02$ and a cliff at $c_w = 0$, suggesting a minimum
threshold. But FI-004 used fixed weights---the Steersman's adaptive decay
was not tested in the low-$c_w$ regime.

CW-001 tests this directly. Configuration: $n = 6$ (RD), identity
initialization, initial $c_w = 0.02$ with the Steersman's adaptive
decay active. Control Law~2 reduces $c_w$ when the co/cross ratio
exceeds the directionality threshold, with a floor at
$0.25 \times c_{w,\text{base}} = 0.005$. The experiment runs for
10{,}000~steps on TinyLlama~1.1B.

\subsubsection{Three-Phase Trajectory}

The CW-001 trajectory reveals three distinct dynamical regimes within
a single training run:

\begin{table}[h]
\centering
\small
\begin{tabular}{@{}rrrrl@{}}
\toprule
Step & Co/Cross & Fiedler & $c_w$ & Phase \\
\midrule
100 & 20:1 & 0.0073 & 0.020 & 1: Initial growth \\
500 & 1{,}592:1 & 0.0004 & 0.015 & 1: Rapid BD formation \\
1{,}500 & 1{,}867:1 & 0.0016 & 0.005 & 2: Plateau onset \\
3{,}300 & 1{,}855:1 & 0.0018 & 0.005 & 2: Sustained plateau \\
4{,}600 & 2{,}252:1 & 0.0019 & 0.005 & 3: Breakout \\
5{,}800 & 3{,}860:1 & 0.0014 & 0.005 & 3: Acceleration \\
7{,}500 & 7{,}884:1 & 0.0007 & 0.005 & 3: Exponential growth \\
8{,}300 & \textbf{15{,}183:1} & \textbf{0.0004} & 0.005 & 3: Peak \\
9{,}000 & 13{,}793:1 & 0.0004 & 0.005 & 4: BD stabilization \\
9{,}500 & 12{,}846:1 & 0.0004 & 0.005 & 4: Oscillation \\
\textbf{10{,}000} & \textbf{13{,}456:1} & \textbf{0.0005} & 0.005 & \textbf{4: Final} \\
\bottomrule
\end{tabular}
\caption{CW-001 trajectory: co/cross ratio, Fiedler eigenvalue, and
contrastive weight across four dynamical phases.}
\label{tab:cw001-trajectory}
\end{table}

\textbf{Phase~1 (steps 0--1{,}500): Rapid initial growth.}
The Steersman decays $c_w$ from 0.020 to the floor of 0.005 within
1{,}500~steps as the co/cross ratio rises rapidly. BD structure forms
quickly during this high-$c_w$ window, reaching ${\sim}1{,}800$:1.

\textbf{Phase~2 (steps 1{,}500--4{,}200): Apparent plateau.}
With $c_w$ floored at 0.005, the co/cross ratio oscillates between
1{,}500:1 and 2{,}000:1 for nearly 3{,}000~steps. The Fiedler
eigenvalue rebounds from its Phase~1 minimum, rising to 0.0020---indicating
that spectral connectivity is \emph{increasing} even as BD strength
appears stalled. This phase was initially interpreted as a terminal
state (adaptive decay defeating whisper-strength). It is not.

\textbf{Phase~3 (steps 4{,}600--8{,}300): Exponential breakout.}
The co/cross ratio breaks above 2{,}000:1 at step~4{,}600 and enters an
exponential growth phase: $2{,}252 \to 3{,}860 \to 7{,}884 \to 12{,}921
\to 15{,}183$:1 (peak at step~8{,}300). The Fiedler eigenvalue declines
monotonically from 0.0019 to 0.0004, entering the same regime as the
aggressive-$c_w$ experiments. The growth rate \emph{accelerates} with
accumulated BD strength, suggesting a positive feedback mechanism: once
block-diagonal structure exceeds a critical threshold (${\sim}2{,}000$:1),
the contrastive loss becomes more effective at reinforcing the existing
partition, creating a self-amplifying cycle.

\textbf{Phase~4 (steps 8{,}300--10{,}000): BD regime stabilization.}
The co/cross ratio oscillates between 12{,}145:1 and 15{,}183:1
(mean~13{,}288:1) without a clear growth trend. The Fiedler eigenvalue
stabilizes at 0.00042--0.00048. Validation loss continues declining
(0.4034 $\to$ 0.4017), indicating that the \emph{language model} is still
improving even as the \emph{bridge structure} has reached equilibrium.
This decoupling---structural stasis while perplexity improves---suggests
the bridge has settled into an attractor from which the remaining
training budget refines weights within the established topology rather
than reshaping it.

\subsubsection{Speed, Destination, or Both?}

At step~10{,}000, CW-001's Fiedler eigenvalue (0.00046) is within one
order of magnitude of H-ch6's final Fiedler (0.00009). Both are deep
in the BD regime. But the co/cross ratios tell a more nuanced story:

\begin{table}[h]
\centering
\small
\begin{tabular}{@{}lrrr@{}}
\toprule
Metric & H-ch6 ($c_w = 0.1$ adaptive) & CW-001 ($c_w = 0.005$ floor) & Ratio \\
\midrule
Co/cross & 70{,}404:1 & 13{,}456:1 & $5.2\times$ \\
Fiedler & 0.00009 & 0.00046 & $5.1\times$ \\
Val loss & ${\sim}0.40$ & 0.4017 & ${\sim}$same \\
\bottomrule
\end{tabular}
\caption{CW-001 vs.\ H-ch6 at matched step count (10{,}000 steps).}
\label{tab:cw001-comparison}
\end{table}

A $5.2\times$ gap persists at matched step count. CW-001's Phase~4
stabilization---oscillating at 12{,}145--15{,}183:1 for 2{,}000~steps
with no growth trend---raises the question: is this a transient plateau
(like Phase~2, which lasted 2{,}700~steps before breaking out) or a
genuine attractor at this $c_w$ floor?

Two interpretations compete. Under the \textbf{speed hypothesis}, CW-001
simply needs more training steps; the Phase~$2 \to 3$ breakout precedent
shows the system can appear stalled and then resume exponential growth.
Under the \textbf{attractor hypothesis}, the contrastive floor sets not
only the rate but the ceiling of structural formation. The $5.2\times$
co/cross gap and the $5.1\times$ Fiedler gap are proportional, consistent
with a $c_w$-dependent attractor strength.

The 24C-002 experiment (\S\ref{sec:24cell}) resolves this question in
favour of the attractor hypothesis. At $n{=}24$ with adaptive $c_w$
settling at $0.025$, the final co/cross is $5{,}983$:1---$6\times$ below
the fixed-$c_w{=}0.1$ run (24C-001, $35{,}808$:1) at matched step count.
The contrastive weight governs not only speed but the \emph{depth} of BD
formation. Higher $c_w$ produces a deeper attractor.

Both interpretations share the same novel core: \emph{any nonzero $c_w$
produces BD structure}, the breakout mechanism at ${\sim}2{,}000$:1 is a
genuine phase transition, and the three-phase incubation trajectory
has no prior art in the adapter literature.

\subsubsection{The Phase~2 Plateau as Incubation}

The 3{,}000-step plateau (Phase~2) is not a failure state. During this
period, the Fiedler eigenvalue rises from 0.0004 to 0.0020---the bridge
is building spectral connectivity while block-diagonal strength appears
stalled. We hypothesize that Phase~2 represents an \emph{incubation
period} in which the bridge accumulates the spectral prerequisites for
exponential BD growth. The Fiedler rebound creates a connected substrate;
Phase~3's breakout then partitions this substrate into blocks.

This two-stage process (connect, then partition) mirrors the
spectral-attractor-then-bifurcation sequence observed across the full
experimental programme (\S\ref{sec:spectral-attractor}--\ref{sec:negative-controls}):
spectral-only training creates connectivity (the attractor); contrastive
loss then breaks this connectivity into directed blocks. CW-001
reproduces this sequence \emph{within a single run} at low~$c_w$.

\subsubsection{Implications}

\begin{enumerate}
\item \textbf{No minimum $c_w$ threshold for BD.} Any nonzero $c_w$
  eventually produces BD structure---the FI-004 ``cliff'' at $c_w = 0$
  is real, but $c_w = 0.005$ is far above it.
\item \textbf{Adaptive decay is not a bottleneck.} The Steersman's
  adaptive Control Law~2, which reduces $c_w$ when BD is detected, does
  not prevent BD formation---it merely slows it. The floor mechanism
  ensures a minimum contrastive signal is always present.
\item \textbf{Training budget trades against $c_w$.} A practitioner with
  abundant compute can use lower $c_w$ for potentially better spectral
  properties; a practitioner needing fast convergence can use higher
  $c_w$. Whether the slow regime produces qualitatively better adapters
  (lower Fiedler $=$ more connected $=$ better information flow) is an
  open empirical question.
\item \textbf{The breakout threshold (${\sim}2{,}000$:1) is a phase
  transition.} The co/cross ratio of ${\sim}2{,}000$:1 marks the onset
  of self-amplifying BD growth at $c_w = 0.005$. Pending verification
  across different $n$ and topologies, this may be universal.
\item \textbf{BD regime stabilization decouples from language modeling.}
  Validation loss continues declining (0.4034 $\to$ 0.4017) during
  Phase~4 while co/cross oscillates without growth. The bridge structure
  reaches equilibrium before the language model does---the final training
  budget refines representations within established topology.
\end{enumerate}


\subsection{Limitations}
\label{sec:disc-limitations}

Several limitations qualify the present results.

\textbf{Pair-based specifications only.}
All topology specifications tested partition channels into co-axial
pairs. More complex topologies---triplets, cycles, cliques of
size~${>}\,2$, hierarchical nesting---would require extending the
contrastive loss formulation. Whether the Steersman can program
non-pair-based structures is untested.

\textbf{Single task family.}
All experiments use instruction-following fine-tuning on
alpaca-cleaned. Whether the four-regime taxonomy, spectral attractor,
and geometric selectivity transfer to other task types
(classification, summarization, code generation) is unknown.

\textbf{Single model family.}
Topology programming experiments use TinyLlama~1.1B. Paper~3
established scale consistency for RD contrastive at 1.1B and 7B, but
the generalization of 4D programming, geometric selectivity, and the
full four-regime taxonomy to larger models is untested.

\textbf{No formal optimality proof.}
We demonstrate that the Steersman programs a given topology, but we
do not prove that any particular topology is optimal for a downstream
task. The relationship between topology specification and task
performance remains to be characterized.

\textbf{24-cell CL2 logging blind spot.}
The 24-cell experiment (24C-001) completed 10{,}000 steps. A logging bug
in Control Law~2 (co/cross tracking gated to $n{=}6$) silenced real-time
co/cross after step~4{,}300 ($5{,}673$:1). Post-hoc recovery (PC-001)
from all 100 saved bridge checkpoints revealed that co/cross grew to
$35{,}808$:1---6.3$\times$ beyond the last real-time measurement. The
contrastive loss itself was applied correctly throughout (via the
cybernetic training script, which has no $n$-gate); only the metric
logging was affected. The bug has been documented for correction.

\textbf{FI-003 sign persistence preliminary.}
The sign persistence experiment (\S\ref{sec:disc-homeostasis}) has
1{,}200~steps of LM-only data (12~checkpoints). The two-timescale decay
trajectory is clear, but the asymptotic equilibrium co/cross ratio
(expected to converge to ${\sim}1$:1) and the long-term validation
loss impact are not yet determined.

\textbf{No theoretical account of the attractor or selectivity.}
The spectral attractor at Fiedler $\approx 0.09$ is empirically
consistent but lacks a theoretical derivation. Similarly, we observe
that only polytope-derived pair specifications produce block-diagonal
structure, but provide no formal characterization of which
specifications constitute valid geometric programs. Both are open
questions for future theoretical work, possibly through random matrix
theory (attractor) or group-theoretic analysis of the contrastive loss
landscape (selectivity).


% ===================================================================
\section{Conclusion}
\label{sec:conclusion}

The Steersman is a general-purpose topology programmer for neural
network adapter bridges---but a selective one. Given a pair
specification derived from the co-axial structure of a polytope, it
compiles the corresponding block-diagonal partition with coupling ratios
exceeding 40{,}000:1. Given a specification lacking geometric
coherence, it drives the bridge into connectivity collapse. The
selectivity is the finding: it confirms that successful topology
programming encodes genuine geometric information, not an artifact of
the loss formulation.

Three polytope geometries are confirmed. The octahedron ($n{=}4$)
produces $2{+}2$ blocks at 473{,}622:1---the programme's strongest
signal. The rhombic dodecahedron ($n{=}6$) produces $3{+}3$ blocks at
70{,}404:1. The tesseract ($n{=}8$) produces $4{+}4$ blocks at
41{,}564:1, extending topology programming to four dimensions. A
24-cell experiment ($n{=}24$, D4 root polytope) confirms a
$12{+}12$ eigenvalue split with co/cross $35{,}808$:1 (post-hoc
PC-001 recovery) and Fiedler stabilized at $5.6 \times 10^{-4}$
for 2{,}100 steps (fixed $c_w{=}0.1$).
In every
confirmed case, the Steersman's control laws, loss formulation, and
hyperparameters are identical; only the pair specification changes.

The spectral attractor at Fiedler $\approx 0.09$ is the universal
unprogrammed state. Spectral regularization alone produces this
equilibrium across $n \in \{3, 4, 8, 12\}$ with 10.4\% variation, and
the emanation architecture converges to the same attractor through a
different mechanism. Contrastive loss breaks this equilibrium by
redirecting connectivity from isotropic distribution to specified
axes---a $1{,}130\times$ Fiedler drop at $n{=}6$.

Geometric coherence is required. Wrong-labels (random partition)
produces a clean $3{+}3$ eigenvalue split but co/cross
$\approx 0$---structure without direction. Resonance
(number-theoretic pairs) produces a chain-like eigenvalue pattern with
neither blocks nor direction. These two negative controls, together
with the three positive results and the spectral attractor baseline,
define four cleanly separated regimes that classify all observed bridge
behavior.

Bridge initialization independence confirms that the topology is
determined by the pair specification, not the initial conditions:
three maximally different corpus-coupled initializations converge to
identical block-diagonal structure
(50--52{,}000:1, Frobenius distance ${\sim}10^{-4}$), with the
Steersman actively correcting 75\% wrong initial signs to the
canonical pattern. Yet this topology requires continuous maintenance:
removing the Steersman causes signs to collapse to random within 100
steps and co/cross to decay exponentially toward the isotropic
equilibrium. The block-diagonal configuration is a Steersman-maintained
fixed point---accessible from any initialization, but unstable under
the language modelling loss alone.

Steersman annealing reveals that the optimal operating point is not at
full strength. Linearly reducing the contrastive weight from $0.10$ to
zero produces five regimes: maintenance, interference (co/cross floor
at 7:1), exponential growth (peak 18{,}671:1 at $c_w = 0.017$), gradual
decay, and a cliff (76\% drop) at $c_w = 0$. The topology programmer
works best at approximately 2\% of its nominal weight---full strength
operates in the interference regime where contrastive and LM gradients
conflict. Even at zero weight, the bridge retains 2{,}942:1,
two orders of magnitude above the isotropic equilibrium, confirming
that annealing imprints topology that partially persists beyond the
Steersman's removal.

The practical implication is precise: adapter topology is now a
designable property, not an emergent accident. The pair specification is
the program, the Steersman is the compiler, and the bridge is the
substrate. The compiler accepts real geometry and rejects everything
else---and it operates most efficiently at near-whisper intensity.


% ===================================================================
% --- references ---
\bibliographystyle{plainnat}
\bibliography{rhombic-paper4}

\end{document}
